% =====================================================================
%  DevOps Technical White Paper
%  "One Verified Path from Pull Request to Production"
%
%  COMPILE WITH XeLaTeX  (Overleaf: Menu -> Compiler -> XeLaTeX)
%  pdfLaTeX will NOT work: this file loads fontspec to honour the
%  Segoe UI formatting requirement.
%
%  ------------------------------------------------------------------
%  FONT NOTE (read once)
%  ------------------------------------------------------------------
%  The brief asks for Segoe UI. Segoe UI is a Microsoft font and is NOT
%  installed on Overleaf, so the preamble tries four sources in order:
%     1. segoeui.ttf etc. uploaded into this Overleaf project
%     2. a system-installed "Segoe UI" family (true on Windows)
%     3. Lato            (humanist sans, present in TeX Live)
%     4. TeX Gyre Heros  (always present; last resort)
%  To get *literal* Segoe UI on Overleaf, upload these four files from
%  C:\Windows\Fonts into the project root:
%     segoeui.ttf  segoeuib.ttf  segoeuii.ttf  segoeuiz.ttf
%  Nothing else needs to change; branch 1 then wins automatically.
%
%  ------------------------------------------------------------------
%  Student and guide details taken from the supplied draft. A course code
%  was not supplied, so the title page names the course without inventing one.
% =====================================================================

\documentclass[10pt,a4paper]{article}

% ---------- geometry: A4 portrait ----------
\usepackage[a4paper,top=2.2cm,bottom=2.2cm,left=2.2cm,right=2.2cm,footskip=1cm]{geometry}

% ---------- fonts (Segoe UI with a graceful fallback chain) ----------
\usepackage{fontspec}
\newif\ifsegoefound \segoefoundfalse

\IfFontExistsTF{segoeui.ttf}{%
  \setmainfont{segoeui.ttf}[
    BoldFont       = segoeuib.ttf,
    ItalicFont     = segoeuii.ttf,
    BoldItalicFont = segoeuiz.ttf,
    Ligatures      = TeX ]
  \segoefoundtrue
}{%
  \IfFontExistsTF{Segoe UI}{%
    \setmainfont{Segoe UI}[Ligatures=TeX]
    \segoefoundtrue
  }{%
    \IfFontExistsTF{Lato}{%
      \setmainfont{Lato}[Ligatures=TeX]
    }{%
      \setmainfont{texgyreheros-regular.otf}[
        BoldFont       = texgyreheros-bold.otf,
        ItalicFont     = texgyreheros-italic.otf,
        BoldItalicFont = texgyreheros-bolditalic.otf,
        Ligatures      = TeX ]
    }%
  }%
}

\IfFontExistsTF{DejaVu Sans Mono}{%
  \setmonofont{DejaVu Sans Mono}[Scale=0.85]
}{%
  \setmonofont{texgyrecursor-regular.otf}[
    BoldFont   = texgyrecursor-bold.otf,
    ItalicFont = texgyrecursor-italic.otf,
    Scale      = 0.85 ]
}

% ---------- layout, colour, typography ----------
\usepackage{setspace}
\setstretch{1.15}                     % brief: 1.15 or 1.5
\usepackage[dvipsnames,table]{xcolor}
\usepackage{titlesec}
\usepackage{enumitem}
\usepackage{booktabs}
\usepackage{tabularx}
\usepackage{array}
\usepackage{caption}
\usepackage{fancyhdr}
\usepackage{microtype}
\usepackage{xurl}
\usepackage{lastpage}
\usepackage[hidelinks,breaklinks=true]{hyperref}
\usepackage{tikz}
\usetikzlibrary{arrows.meta,positioning,calc}

% long monospace paths must not produce overfull lines
\setlength{\emergencystretch}{3em}
\renewcommand{\refname}{References}

\definecolor{slate}{HTML}{1F2937}
\definecolor{steel}{HTML}{334155}
\definecolor{muted}{HTML}{64748B}
\definecolor{rule}{HTML}{94A3B8}
\definecolor{srcfill}{HTML}{EEF2FF}
\definecolor{cifill}{HTML}{ECFDF5}
\definecolor{relfill}{HTML}{FFF7ED}
\definecolor{depfill}{HTML}{EFF6FF}
\definecolor{verfill}{HTML}{FDF2F8}
\definecolor{alert}{HTML}{B91C1C}
\definecolor{headrow}{HTML}{F1F5F9}

\color{slate}

% ---------- the details to personalise ----------
\newcommand{\mytitle}{One Verified Path from Pull Request to Production}
\newcommand{\mysubtitle}{A Pipeline-First DevOps Architecture for a Git-Driven Online
Assignment Submission Portal}
\newcommand{\myname}{Ayush Gole}
\newcommand{\myreg}{23102A0005}
\newcommand{\mydept}{Department of Computer Engineering, Vidyalankar Institute of Technology}
\newcommand{\myguide}{Dr. Snehal Andhare}
\newcommand{\mycourse}{DevOps}

% ---------- headings: 12 pt bold ----------
\titleformat{\section}{\fontsize{12}{14.4}\selectfont\bfseries\color{steel}}
            {\thesection.}{0.6em}{}
\titlespacing*{\section}{0pt}{9pt plus 2pt minus 1pt}{4pt}
\titleformat{\subsection}{\fontsize{10.5}{12.6}\selectfont\bfseries\color{steel}}
            {\thesubsection}{0.6em}{}
\titlespacing*{\subsection}{0pt}{6pt plus 1pt minus 1pt}{2pt}

\captionsetup{font={small},labelfont={bf},skip=5pt,justification=justified}
\setlist[itemize]{leftmargin=1.25em,itemsep=1.5pt,topsep=2.5pt,parsep=0pt}
\setlist[enumerate]{leftmargin=1.6em,itemsep=2pt,topsep=2.5pt,parsep=0pt}

% ---------- running footer ----------
\pagestyle{fancy}
\fancyhf{}
\renewcommand{\headrulewidth}{0pt}
\fancyfoot[L]{\footnotesize\color{muted}DevOps Technical White Paper -\ \myname}
\fancyfoot[R]{\footnotesize\color{muted}Page \thepage\ of \pageref{LastPage}}

% table helpers
\newcolumntype{L}[1]{>{\raggedright\arraybackslash}p{#1}}
\renewcommand{\arraystretch}{1.12}

\begin{document}

% =====================================================================
% PAGE 1 -- TITLE
% =====================================================================
\thispagestyle{empty}
\begin{center}
\vspace*{1.6cm}

{\footnotesize\color{muted}INDIVIDUAL TECHNICAL WHITE PAPER}\\[0.25cm]
{\color{rule}\rule{0.42\textwidth}{0.6pt}}\\[0.9cm]

{\fontsize{16}{19.2}\selectfont\bfseries \mytitle\par}
\vspace{0.45cm}
{\fontsize{12}{15}\selectfont\bfseries\color{steel} \mysubtitle\par}

\vspace{0.9cm}
{\color{rule}\rule{0.42\textwidth}{0.6pt}}\\[1.1cm]

{\fontsize{11}{13.2}\selectfont
\textbf{Submitted by}\\[0.2cm]
{\fontsize{12}{14.4}\selectfont\bfseries \myname}\\[0.15cm]
\myreg\\[0.15cm]
\mydept\par}

\vspace{0.9cm}

{\fontsize{11}{13.2}\selectfont
\textbf{Guide}\\[0.2cm]
{\fontsize{12}{14.4}\selectfont\bfseries \myguide}\par}

\vspace{0.9cm}

{\fontsize{11}{13.2}\selectfont
\textbf{Course}\\[0.2cm]
\mycourse\par}

\vfill

{\footnotesize\color{muted}
Subject of the paper: the delivery pipeline of \texttt{online-assignment-portal} -\\
a Spring Boot modular monolith delivered by GitHub, Jenkins, Maven, Testcontainers,\\
Selenium, Docker, a container registry, and Ansible.\\[0.5cm]
October 2026}

\vspace{0.4cm}
\end{center}

\newpage

% =====================================================================
% PAGE 2 -- ABSTRACT, INTRODUCTION, PROBLEM STATEMENT
% =====================================================================
\section*{Abstract}

A coursework submission portal is a small application; delivering it dependably is not a small
problem. This paper takes the delivery pipeline of a Git-driven online assignment submission
portal -\ a Spring Boot modular monolith in which students upload PDF coursework and
reviewers record one terminal decision -\ and proposes an architecture in which a
change reaches a running host along exactly one automated path. The governing commitment is that
\emph{the artefact which is tested is the artefact which runs}: Maven produces one executable
JAR, a multi-stage \texttt{Dockerfile} packages that JAR without rebuilding it, a Git tag of the
form \texttt{vX.Y.Z} is the only thing that names a version, and the deployed container
republishes that version at \texttt{/actuator/info} so a post-deployment smoke test can prove
that the release deployed is the release answering. Integration tests run against a containerised
MySQL~8.4 rather than an in-memory substitute, so migrations and constraints are exercised by the
engine that runs in production; provisioning is expressed as Ansible roles that leave the stack
under \texttt{systemd}; and recovery is a designed path built from a pre-release gate, an
append-only release ledger, and a rollback to the last version a ledger row marks stable. The
paper compares this design with manual deployment, hosted CI and full orchestration, separates
what has been executed from what is written but unproven, and sets out an extension path covering
TLS, supply-chain attestation, zero-downtime release and DORA instrumentation.

\section{Introduction}

The application is deliberately modest: students authenticate, upload a PDF and read the outcome;
reviewers see a queue of submitted work and record one decision that is final. The real
deliverable is the path that carries a change -\ a pull request on GitHub, review,
Jenkins, Maven, unit tests, integration tests, Selenium end-to-end tests, a Docker image, a
registry, Ansible provisioning, deployment, a health check and a rollback -\ repeatably
and with evidence. Complexity added to the application is therefore a cost charged against the
pipeline budget, and a feature the pipeline cannot test, build and deploy is not finished.

That modest scope still contains what makes deployment hard in practice: \textbf{state, in two
incompatible forms}. A submission is a row in MySQL \emph{and} a file of bytes on a volume, and
the two cannot be written in one transaction. A container filesystem is ephemeral, so the bytes
must outlive the container that wrote them; a migration moves forward while a rollback moves
backward; and an upload is hostile input until proven otherwise. A pipeline here must therefore
replace a process \emph{while} preserving two kinds of durable state and remaining able to go
back. This paper contributes an analysis of why common delivery approaches under-serve a
single-host, state-bearing deployment of this size (Section~3); a pipeline architecture whose
distinguishing rules are \emph{tag-as-version} and \emph{verification-by-version-assertion}
(Section~4); an implementation path mapped onto concrete artefacts rather than generic advice
(Section~5); and an account that separates executed stages from written-but-unproven ones
(Section~6).

\section{Problem Statement}

Six defects recur in student and small-team delivery, and each is a requirement in disguise.

\begin{enumerate}[label=\textbf{D\arabic*.},leftmargin=2.3em]
  \item \textbf{Artefact drift.} The software is rebuilt per environment, so the binary that was
        tested is not the binary that runs; every rebuild is an untested change, and
        configuration baked at build time makes promotion impossible.
  \item \textbf{Substituted test dependencies.} Persistence is tested on an in-memory database, so
        migrations, \texttt{CHECK} constraints, unique indexes and real locking behaviour are
        never exercised where failure is cheap.
  \item \textbf{Unidentifiable releases.} A mutable \texttt{latest} tag means no host can answer
        ``which version is running here'', and two machines pulling one tag can run different
        code.
  \item \textbf{Recovery by improvisation.} Rollback means rebuilding an older commit under
        pressure, and the identity of the last version that actually worked is recorded nowhere a
        script can read.
  \item \textbf{State asymmetry.} Rows and bytes are protected by different mechanisms and
        restored by different commands, so a backup of one is mistaken for a backup of the
        submission.
  \item \textbf{Verification by anecdote, secrets by convenience.} A deployment is called
        successful because a page loaded rather than because the build asserted its own version,
        while credentials sit in committed files because externalising them was never the easy
        path.
\end{enumerate}

\noindent\textbf{Objective.} Build one artefact once, prove it against production-grade
dependencies, publish it under an immutable version, deploy it by script to a provisioned host,
verify it by asserting that the running version equals the released version, and keep it
reversible to a recorded good version -\ with no secret in the repository at any point.

\newpage

% =====================================================================
% PAGE 3 -- EXISTING APPROACHES, TOOLS AND LIMITATIONS
% =====================================================================
\section{Existing Approaches, Tools and Limitations}

\subsection{Manual and script-driven release}
The cheapest approach copies a JAR to a host over \texttt{scp} and restarts a service. It needs
no infrastructure and is the baseline most projects start from, but it supplies none of D1--D6:
there is no gate, no recorded identity, no reproducible host state, and the knowledge of how to
deploy lives in one person's shell history. The systematic review by Shahin \textit{et al.}
records exactly this pattern -\ manual, undocumented deployment -\ as a
leading obstacle to continuous delivery adoption~\cite{shahin2017,chen2015}, and Humble and
Farley treat the automated, repeatable deployment as the precondition for everything
else~\cite{humble2010}.

\subsection{CI platforms: hosted versus self-hosted}
Hosted runners (GitHub Actions, GitLab CI/CD, CircleCI) remove agent maintenance and express the
pipeline as versioned YAML, which is the correct instinct: pipeline-as-code. Their limitations
here are contextual rather than technical -\ build minutes are metered, the Docker
daemon is someone else's, and the course objective is an owned, inspectable controller.
Jenkins with a declarative \texttt{Jenkinsfile} keeps the pipeline definition in the repository
while leaving the agent and its Docker daemon under local control~\cite{jenkins}. The price is
real: plugin surface area, agent upkeep, and credential hygiene become the project's own
responsibility, which is why credentials are referenced by identifier and bound only for the
length of a single stage.

\subsection{Test infrastructure}
H2 or HSQLDB make persistence tests fast, and that speed is bought by testing a different
database. Testcontainers inverts the trade: it starts a real MySQL~8.4 in Docker for the duration
of the suite, so migrations and constraints are proven against the production
engine~\cite{testcontainers}. Selenium proves that a human can complete a journey through a real
browser~\cite{selenium}, but it is minutes-slow and becomes flaky when asked to prove business
rules that a unit test could settle. The practical rule is to test at the lowest layer that can
prove the behaviour and let the pyramid stay a pyramid.

\subsection{Packaging, registry and orchestration}
A fat JAR on a host leaves the runtime, the JVM flags, and the filesystem layout implicit. A
container image makes them explicit, and multi-stage builds keep the published image free of
build tooling~\cite{dockermulti}. Kubernetes with Helm and a GitOps reconciler such as Argo CD is
the industry default for multi-instance services, but it imposes a control plane that a
single-host, single-instance MVP cannot repay; Docker Compose on an Ansible-provisioned host is
the proportionate choice at this scale~\cite{ansible}. Two registry habits defeat identity
regardless of orchestrator: mutable tags, and the absence of provenance metadata such as an SBOM
or a signature~\cite{ociimage,slsa}.

\begin{table}[h]
\centering
\caption{Candidate approaches assessed against the objectives of Section~2.}
\vspace{-2pt}
{\footnotesize
\begin{tabularx}{\textwidth}{L{3.05cm} L{3.35cm} X}
\toprule
\rowcolor{headrow}\textbf{Approach} & \textbf{Representative tools} &
\textbf{Principal limitation in this context} \\
\midrule
Manual / script release & \texttt{scp}, \texttt{systemctl} & No gate, no identity, no recorded
way back; host state is undocumented (D1, D3, D4). \\
Hosted CI & GitHub Actions, GitLab CI & Metered minutes and a borrowed Docker daemon; the
controller is not ours to inspect or extend. \\
Self-hosted CI & Jenkins declarative pipeline & Plugin and agent upkeep; credential hygiene
becomes a project responsibility. \\
In-memory test DB & H2, HSQLDB & Different dialect: migrations and \texttt{CHECK} constraints
are never exercised (D2). \\
Container orchestration & Kubernetes, Helm, Argo CD & Control-plane cost is disproportionate to
one instance; a second system to provision and secure. \\
Registry defaults & \texttt{latest}, unsigned images & Version cannot be asserted at runtime; no
provenance for what was shipped (D3). \\
\bottomrule
\end{tabularx}}
\end{table}

\subsection{The gap this paper addresses}
Each tool above is sound in its own domain, and none of them, adopted by default, delivers the
four properties together at this scale: one artefact promoted rather than rebuilt; test fidelity
against the production database engine; a version identity that the running system can assert
about itself; and a scripted, evidence-recording way back. The proposal that follows composes
established tools so that these four hold simultaneously, and deliberately declines the
orchestration tier whose cost the application cannot yet justify.

\newpage

% =====================================================================
% PAGE 4 -- PROPOSED TECHNICAL SOLUTION AND ARCHITECTURE
% =====================================================================
\section{Proposed Technical Solution and Architecture}

\subsection{Design rules}
\begin{enumerate}[label=\textbf{R\arabic*.},leftmargin=2.3em]
  \item \textbf{One artefact, promoted not rebuilt.} \texttt{./mvnw} produces the JAR; the
        \texttt{Dockerfile} runs no Maven and only packages what the build already produced.
  \item \textbf{The tag is the version.} A \texttt{vX.Y.Z} Git tag is the sole source of a
        version~\cite{semver}; a feature branch reaches neither registry nor host, and
        \texttt{latest} is refused by the publish step rather than merely discouraged.
  \item \textbf{Prove it at the lowest layer that can.} Unit tests for rules, Testcontainers
        integration tests for persistence, security wiring and migrations, Selenium only for the
        browser journey.
  \item \textbf{Configuration is injected, never baked.} Defaults in \texttt{application.yml},
        non-secret profile values in profile files, every secret from the environment or a
        managed store~\cite{owasp2021}.
  \item \textbf{Recovery is a designed path.} A pre-release gate, an append-only ledger naming
        the last stable version, and a rollback that is the same code path whether a human or a
        failing stage invokes it.
  \item \textbf{A claim is worth its transcript.} Every stage emits an artefact -\ test
        reports, screenshots, a database dump, a health and smoke transcript, a ledger row
        -\ so that ``it works'' is a record rather than a recollection.
\end{enumerate}

\subsection{Architecture}

\begin{figure}[h]
\centering
\resizebox{\textwidth}{!}{%
\begin{tikzpicture}[
  x=1cm, y=1cm,
  box/.style={draw=rule, line width=0.45pt, rounded corners=2.5pt, align=center,
              text width=2.8cm, minimum height=1.15cm, inner sep=3.5pt, font=\scriptsize},
  src/.style={box, fill=srcfill}, ci/.style={box, fill=cifill},
  rel/.style={box, fill=relfill}, dep/.style={box, fill=depfill},
  ver/.style={box, fill=verfill},
  lane/.style={rotate=90, font=\scriptsize\bfseries, text=muted},
  flow/.style={-{Stealth[length=4.5pt,width=3.5pt]}, draw=steel, line width=0.5pt},
  back/.style={-{Stealth[length=4.5pt,width=3.5pt]}, draw=alert, line width=0.5pt, dashed},
  tag/.style={font=\scriptsize\itshape, text=muted, inner sep=1.5pt}
]

% ---- lane labels ----
\node[lane] at (-1.0,0)    {SOURCE};
\node[lane] at (-1.0,-2.1) {BUILD};
\node[lane] at (-1.0,-4.2) {RELEASE};
\node[lane] at (-1.0,-6.3) {DEPLOY};
\node[lane] at (-1.0,-8.4) {VERIFY};

% ---- row A: source and review ----
\node[src] (a1) at (1.6,0)  {\textbf{Developer}\\\texttt{./mvnw clean verify}\\locally, before pushing};
\node[src] (a2) at (5.1,0)  {\textbf{GitHub}\\\texttt{feature/*} branch\\pull request, protected \texttt{main}};
\node[src] (a3) at (8.6,0)  {\textbf{Peer review}\\PR template: problem,\\approach, verification};
\node[src] (a4) at (12.1,0) {\textbf{Jenkins trigger}\\declarative \texttt{Jenkinsfile}\\one build per branch};
\node[src] (a5) at (15.6,0) {\textbf{Secrets}\\credentials store (release)\\Vault (provision)\\host \texttt{.env} (runtime)\\\textit{bound per stage, never committed}};

% ---- row B: CI ----
\node[ci] (b1) at (1.6,-2.1)  {\textbf{Validate + Build}\\Maven wrapper, Java 21\\enforcer rules};
\node[ci] (b2) at (5.1,-2.1)  {\textbf{Unit tests}\\JUnit 5, no Spring\\context (Surefire)};
\node[ci] (b3) at (8.6,-2.1)  {\textbf{Integration tests}\\Testcontainers MySQL 8.4\\+ Flyway V1--V3 (Failsafe)};
\node[ci] (b4) at (12.1,-2.1) {\textbf{Package}\\one executable JAR\\\textit{built once}};
\node[ci] (b5) at (15.6,-2.1) {\textbf{Selenium E2E}\\browser journeys vs\\an ephemeral stack};

% ---- row C: release ----
\node[rel] (c1) at (1.6,-4.2)  {\textbf{Release Version}\\git tag \texttt{vX.Y.Z}\\\textit{the only version source}};
\node[rel] (c2) at (5.1,-4.2)  {\textbf{Docker Build}\\multi-stage, layered JAR\\non-root UID 1001};
\node[rel] (c3) at (8.6,-4.2)  {\textbf{Version baked in}\\\texttt{APP\_VERSION} build arg\\$\rightarrow$ \texttt{/actuator/info}};
\node[rel] (c4) at (12.1,-4.2) {\textbf{Publish Image}\\Docker Hub\\\texttt{acct/repo:vX.Y.Z}};
\node[rel] (c5) at (15.6,-4.2) {\textbf{Registry}\\immutable tag per release\\\texttt{latest} refused};

% ---- row D: provision and deploy ----
\node[dep] (d1) at (1.6,-6.3)  {\textbf{Ansible} \texttt{site.yml}\\\texttt{common} \textperiodcentered\ \texttt{docker} \textperiodcentered\ \texttt{application}};
\node[dep] (d2) at (5.1,-6.3)  {\textbf{Host} \texttt{/opt/assignment-portal}\\\texttt{compose.prod.yaml} + \texttt{.env}\\\texttt{systemd} unit};
\node[dep] (d3) at (8.6,-6.3)  {\textbf{Pre-release gate}\\volumes verified \textperiodcentered\ MySQL dump\\\texttt{releases/<v>.env} written};
\node[dep] (d4) at (12.1,-6.3) {\textbf{Deploy}\\\texttt{compose pull} then\\\texttt{up -{}-no-build -{}-wait}};
\node[dep] (d5) at (15.6,-6.3) {\textbf{Running stack}\\app :8080 + MySQL 8.4\\volumes: uploads, data};

% ---- row E: verify and recover ----
\node[ver] (e1) at (1.6,-8.4)  {\textbf{Health gate}\\\texttt{GET /actuator/health}\\polled from the agent};
\node[ver] (e2) at (5.1,-8.4)  {\textbf{Smoke test}\\\texttt{/actuator/info} version\\$=$ the released tag};
\node[ver] (e3) at (8.6,-8.4)  {\textbf{Release ledger}\\\texttt{releases/ledger.tsv}\\append-only outcome};
\node[ver] (e4) at (12.1,-8.4) {\textbf{Rollback}\\last version a \texttt{success}\\row marks stable};
\node[ver] (e5) at (15.6,-8.4) {\textbf{Backup / restore}\\DB dump + upload bytes\\\textit{both halves of a row}};

% ---- in-row flow ----
\foreach \a/\b in {a1/a2, a2/a3, a3/a4, b1/b2, b2/b3, b3/b4, b4/b5,
                   c1/c2, c2/c3, c3/c4, c4/c5, d1/d2, d2/d3, d3/d4, d4/d5,
                   e1/e2, e2/e3, e3/e4, e4/e5}
  \draw[flow] (\a.east) -- (\b.west);

% ---- row-to-row flow ----
\draw[flow] (a4.south) -- ++(0,-0.46) -| (b1.north);
\draw[flow] (b5.south) -- ++(0,-0.46) -| (c1.north);
\draw[flow] (c5.south) -- ++(0,-0.46) -| (d1.north);
\draw[flow] (d5.south) -- ++(0,-0.46) -| (e1.north);

% ---- recovery loop: the designed way back ----
\draw[back] (e4.north) -- (d4.south);

% ---- failure annotation ----
\node[tag, anchor=west] at (8.3,-9.75) {health or smoke failure $\Rightarrow$ ledger
records \texttt{failure}, rollback fires};
\end{tikzpicture}}
\caption{The proposed pipeline. A change moves down the lanes exactly once: the JAR is built in
the BUILD lane and never rebuilt, the RELEASE lane bakes the Git tag into the image and
publishes it under that immutable tag, the DEPLOY lane pulls that tag onto a provisioned host
behind a pre-release gate, and the VERIFY lane proves the running version equals the released
version before recording the outcome. The dashed path is the designed way back. Secrets (top
right) are bound into a single stage by identifier -\ the registry login in the release
lane, the SSH key and host \texttt{.env} in the deploy lane -\ and appear in no
committed file.}
\end{figure}

\newpage

% =====================================================================
% PAGE 5 -- IMPLEMENTATION STEPS
% =====================================================================
\section{Implementation Steps}

\begin{enumerate}[leftmargin=1.9em]
  \item \textbf{Govern the source.} \texttt{main} is protected and accepts no direct push. Work
        arrives on a \texttt{feature/}, \texttt{fix/}, \texttt{docs/} or \texttt{refactor/} branch
        as a pull request whose template states the problem, the approach and the verification
        actually performed. Commits follow Conventional Commits~\cite{convcommits}, so history is
        machine-readable and release notes derive from it rather than from memory; branches are
        kept short-lived so integration stays continuous~\cite{fowlerci}.
  \item \textbf{Express the pipeline as code.} A declarative \texttt{Jenkinsfile} in the
        repository defines every stage~\cite{jenkins}. Options pin the behaviour that otherwise
        bites: timestamped logs, a 90-minute timeout, concurrent builds disabled so two runs
        cannot race for the same database port and upload volume, and a bounded build history.
        Credentials appear only as identifiers; values are bound into one stage and nowhere else.
  \item \textbf{Build with a pinned toolchain.} The Maven Wrapper is the canonical build
        mechanism, so the agent needs nothing on \texttt{PATH} and every machine resolves the same
        Maven version~\cite{maven}. Enforcer rules fail the build on a banned dependency or a
        wrong Java release rather than at runtime.
  \item \textbf{Run the unit suite first.} JUnit~5 tests with no Spring context assert the
        business rules in milliseconds -\ decision terminality, the required rejection
        reason, PDF and size validation -\ so the fast signal arrives before anything
        expensive starts.
  \item \textbf{Integration-test against the production engine.} Failsafe tests start MySQL~8.4
        through Testcontainers and Flyway applies \texttt{V1}--\texttt{V3} into that
        container~\cite{flyway}, so the suite proves what only a real engine can: migrations,
        \texttt{CHECK} constraints, unique indexes, optimistic locking under two concurrent
        reviewers, deny-by-default routes, CSRF enforcement, and hostile uploads -\
        oversized files, non-PDF bytes behind a \texttt{.pdf} name, and path-traversal filenames.
  \item \textbf{Package once.} \texttt{./mvnw package} produces the single executable JAR that
        every later stage consumes; no stage after this one compiles source.
  \item \textbf{Prove the browser journey against a real stack.} Scripted helpers start the
        database, start the application, wait for \texttt{/actuator/health}, run the Selenium
        suite, collect diagnostics and tear the environment down; reports and failure screenshots
        are archived with the build~\cite{selenium}.
  \item \textbf{Derive the version from the tag.} The \texttt{vX.Y.Z} tag on the built commit
        \emph{is} the version~\cite{semver}. The optional \texttt{RELEASE\_VERSION} parameter is an
        assertion, not an input: when it disagrees with the tag the build fails. A build without a
        tag publishes nothing.
  \item \textbf{Build the image deterministically.} A multi-stage \texttt{Dockerfile} extracts the
        JAR into Spring Boot layers so dependency layers cache across releases, then assembles an
        \texttt{eclipse-temurin:21-jre} runtime that runs as a non-root user (UID/GID~1001),
        declares a \texttt{HEALTHCHECK}, bakes the version in through an \texttt{APP\_VERSION}
        build argument, bounds heap by container memory and stops on
        \texttt{SIGTERM}~\cite{dockermulti}. The build context admits exactly one file, the JAR.
  \item \textbf{Publish under an immutable tag.} \texttt{docker login} receives the token on
        standard input, never as an argument; the image is pushed as
        \texttt{<account>/<repo>:vX.Y.Z} and a request to publish \texttt{latest} is refused. The
        pushed digest is the artefact identity downstream~\cite{ociimage}.
  \item \textbf{Provision the host with Ansible.} \texttt{site.yml} runs three roles:
        \texttt{common} (base packages, time sync), \texttt{docker} (engine, Compose~v2 plugin,
        Python SDK, and the two \texttt{daemon.json} settings this project needs -\
        bounded container logs and \texttt{live-restore}) and \texttt{application} (deploy
        directory, composition, \texttt{.env}, named volumes, image, and a \texttt{systemd} unit
        so the stack is operated as a service)~\cite{ansible}. Secrets live in an
        \texttt{ansible-vault} file referenced explicitly from the playbook, and the role copies
        the \emph{same} composition the deployment path uses, so two files cannot drift.
  \item \textbf{Deploy behind a gate.} The agent copies the host scripts over SSH and runs the
        pre-release gate: volumes verified, a MySQL dump taken, and the release's metadata written
        under \texttt{releases/}; a host that has released before and has lost a volume is refused
        rather than upgraded. The gate rewrites exactly two non-secret keys -\ image tag
        and environment name -\ then \texttt{compose pull} and
        \texttt{up -{}-no-build -{}-wait} swap in the published container. No \texttt{build:}
        exists in the deployment composition, so a host runs the tested artefact or fails.
  \item \textbf{Verify, record, and be able to return.} The agent polls
        \texttt{/actuator/health}~\cite{springboot}, then the smoke test asserts that
        \texttt{/actuator/info} reports the version just released and that a protected route is
        still guarded -\ which is what separates ``up'' from ``correct''. The outcome is
        appended to the release ledger, and on failure the stage invokes the same rollback command
        an operator would run by hand, restoring the last version a \texttt{success} row marks
        stable. One code path, exercised by both routes.
\end{enumerate}

\newpage

% =====================================================================
% PAGE 6 -- OUTPUT, BENEFITS, FUTURE SCOPE, CONCLUSION, REFERENCES
% =====================================================================
\section{Output and Verification Status}

\textbf{Expected pipeline outputs.} Once run, the stages are designed to retain Surefire and
Failsafe reports, Selenium screenshots and diagnostics, the packaged JAR, an image tagged
\texttt{vX.Y.Z} with its registry digest, a pre-release database dump, a
\texttt{releases/<version>.env} context file, a health and smoke transcript, and a release-ledger
row. These are expected outputs of the design, not artefacts currently captured in the repository.

\textbf{What has been executed.} Three local \texttt{clean verify} runs were recorded on
2026-10-02; none established a green full suite. In run 1, both source sets compiled and all 60
unit tests passed, but the 112 integration tests could not start because Docker was unavailable.
With Docker available, run 2 executed 112 integration tests and found two defects: a Thymeleaf
fragment parameter shadowed the upload title model attribute, and the error-page test did not send
the HTML \texttt{Accept} header. Run 3 executed 113 integration tests, confirmed those two fixes,
and exposed a third defect: ordinary HTML comments were reaching rendered pages, including notes
about internal classes and configuration. The repository records fixes for these defects, but no
follow-up run has established that the complete suite is green. The cases provide concrete
evidence that testing found defects, while the result remains explicitly unresolved.

\textbf{What has not been executed.} The Jenkins pipeline has not run on an agent. No Docker image
has been built by the pipeline, pushed to a registry, or deployed; no Ansible playbook or host
release has run, and no rollback drill has been performed. Selenium and delivery outcomes remain
unverified as documented by the project. The implementation is therefore a proposed, documented
delivery path whose operational claims still require captured runs under
\texttt{docs/evidence/}.

\section{Benefits}

\begin{itemize}
  \item \textbf{No drift (D1).} One JAR, one image, one digest promoted across environments;
        configuration arrives at deploy time, so the tested artefact is the running artefact.
  \item \textbf{Defects caught where they are cheap (D2).} Migration, constraint and concurrency
        faults surface pre-merge against MySQL~8.4, not during a demonstration.
  \item \textbf{Answerable identity (D3).} ``Which version is running?'' is an HTTP request;
        ``which version was last good?'' is a ledger row.
  \item \textbf{Bounded recovery (D4, D5).} Rollback is one command on a recorded target, the
        automatic and manual paths are the same code, and both halves of a submission -\
        row and bytes -\ have a named backup and restore procedure.
  \item \textbf{Security by default (D6).} No secret in the repository, non-root container,
        externalised configuration, CSRF enabled, server-side authorisation on every request, and
        uploads served only as attachments~\cite{nist800190,owasp2021}.
  \item \textbf{Evidence as a by-product.} Lead time, change-failure rate and restore time are
        derivable from build records and the ledger rather than estimated -\ the
        measurement discipline the DORA research associates with higher delivery
        performance~\cite{forsgren2018,dora2024}.
\end{itemize}

\section{Future Scope}

\begin{itemize}
  \item \textbf{Transport and perimeter.} A TLS-terminating reverse proxy with automated
        certificate renewal, HSTS, and a host firewall, which together allow the session cookie to
        become \texttt{Secure}.
  \item \textbf{Supply-chain assurance.} An SBOM per image, signing with Sigstore/cosign, SLSA
        provenance attestation, digest-pinned base images, and a vulnerability gate that fails the
        build on fixable critical findings~\cite{slsa,nist800190}.
  \item \textbf{Zero-downtime release.} Blue-green on a single host via two Compose projects
        behind the proxy~\cite{fowlerbg}, with expand-and-contract migrations so a rollback never
        meets a schema it cannot read -\ the one gap that makes the current rollback a
        last resort rather than a routine.
  \item \textbf{Rehearsed recovery.} Backup retention, an off-host copy, and a scheduled restore
        drill run as a pipeline job; until a restore has been performed, a backup is a hypothesis.
  \item \textbf{Observability and metrics.} OpenTelemetry traces and Prometheus metrics behind
        the existing request-identifier correlation seam, plus explicit DORA
        instrumentation~\cite{dora2024}.
  \item \textbf{Policy and architecture enforcement.} Static policy checks on the composition and
        playbooks, and dependency-direction tests promoting the package rules from prose into a
        failing build.
  \item \textbf{A justified scale-out path.} Should multiple instances become warranted,
        object storage replaces the upload volume and an orchestrator replaces Compose
        -\ a change that supersedes recorded architecture decisions rather than
        silently diverging from them.
\end{itemize}

\section{Conclusion}

The application in this project is small by design; the delivery path is the engineering. By
insisting that one artefact is built once, proven against the engine it will actually run on,
named by an immutable tag it can assert about itself, deployed behind a gate that protects both
halves of its state, and reversible to a version a ledger records as stable, the pipeline turns
deployment from an event that is survived into a procedure that is repeated. The design is
complete and documented; the work that remains is the work this paper refuses to skip over
-\ executing it, and keeping the transcript.

\section{References}
{\fontsize{8.6}{10.3}\selectfont
\begin{thebibliography}{99}
\setlength{\itemsep}{1.5pt}\setlength{\parskip}{0pt}

\bibitem{humble2010} J. Humble and D. Farley, \textit{Continuous Delivery: Reliable Software
Releases through Build, Test, and Deployment Automation}. Boston, MA, USA: Addison-Wesley, 2010.

\bibitem{forsgren2018} N. Forsgren, J. Humble, and G. Kim, \textit{Accelerate: The Science of
Lean Software and DevOps}. Portland, OR, USA: IT Revolution Press, 2018.

\bibitem{shahin2017} M. Shahin, M. A. Babar, and L. Zhu, ``Continuous integration, delivery and
deployment: A systematic review on approaches, tools, challenges and practices,''
\textit{IEEE Access}, vol. 5, pp. 3909--3943, 2017.

\bibitem{chen2015} L. Chen, ``Continuous delivery: Huge benefits, but challenges too,''
\textit{IEEE Software}, vol. 32, no. 2, pp. 50--54, Mar./Apr. 2015.

\bibitem{fowlerci} M. Fowler, ``Continuous integration,'' martinfowler.com. [Online]. Available:
\url{https://martinfowler.com/articles/continuousIntegration.html}

\bibitem{fowlerbg} M. Fowler, ``BlueGreenDeployment,'' martinfowler.com. [Online]. Available:
\url{https://martinfowler.com/bliki/BlueGreenDeployment.html}

\bibitem{dora2024} DORA, ``Accelerate State of DevOps Report,'' Google Cloud / DORA research
program. [Online]. Available: \url{https://dora.dev/research/}

\bibitem{jenkins} Jenkins Project, ``Pipeline -\ Jenkins user documentation.''
[Online]. Available: \url{https://www.jenkins.io/doc/book/pipeline/}

\bibitem{maven} Apache Maven Project, ``Introduction to the build lifecycle.'' [Online].
Available: \url{https://maven.apache.org/guides/introduction/introduction-to-the-lifecycle.html}

\bibitem{testcontainers} Testcontainers, ``Testcontainers for Java -\ documentation.''
[Online]. Available: \url{https://java.testcontainers.org/}

\bibitem{selenium} Selenium Project, ``WebDriver -\ Selenium documentation.''
[Online]. Available: \url{https://www.selenium.dev/documentation/webdriver/}

\bibitem{dockermulti} Docker Inc., ``Multi-stage builds -\ Docker Docs.'' [Online].
Available: \url{https://docs.docker.com/build/building/multi-stage/}

\bibitem{ociimage} Open Container Initiative, ``OCI image format specification.'' [Online].
Available: \url{https://github.com/opencontainers/image-spec}

\bibitem{ansible} Red Hat, ``Ansible documentation -\ tips and tricks.'' [Online].
Available: \url{https://docs.ansible.com/ansible/latest/tips_tricks/index.html}

\bibitem{flyway} Redgate, ``Flyway documentation -\ migrations.'' [Online]. Available:
\url{https://documentation.red-gate.com/flyway}

\bibitem{springboot} Spring, ``Spring Boot reference documentation -\ production-ready
features (Actuator).'' [Online]. Available:
\url{https://docs.spring.io/spring-boot/reference/actuator/index.html}

\bibitem{nist800190} M. Souppaya, J. Morello, and K. Scarfone, ``Application container security
guide,'' National Institute of Standards and Technology, Gaithersburg, MD, USA,
NIST SP 800-190, Sep. 2017.

\bibitem{owasp2021} OWASP Foundation, ``OWASP Top 10:2021.'' [Online]. Available:
\url{https://owasp.org/Top10/}

\bibitem{semver} T. Preston-Werner, ``Semantic versioning 2.0.0.'' [Online]. Available:
\url{https://semver.org/}

\bibitem{convcommits} Conventional Commits, ``Conventional Commits 1.0.0.'' [Online]. Available:
\url{https://www.conventionalcommits.org/en/v1.0.0/}

\bibitem{slsa} Open Source Security Foundation, ``SLSA: Supply-chain Levels for Software
Artifacts, specification v1.0.'' [Online]. Available: \url{https://slsa.dev/spec/v1.0/}

\end{thebibliography}}

\end{document}
